% ============================================================================
% CHAPITRE 1 — CONTEXTE, \'ETAT DE L'ART ET PROBL\'eMATIQUE
% Section 1.2 : Contexte
% source : partie1/contexte.tex
% ============================================================================

\section{CONTEXTE}
\label{sec:contexte}
\begin{sloppypar}
La cryptographie et le droit ont toujours entretenu une relation asym\'etrique~\cite{Mason07,Twining1999}.
La premi\`ere produit des garanties absolues, prouvables sous hypoth\`eses
math\'ematiques~\cite{Goldreich01,GMR88}, alors que le second exige des preuves que des
acteurs non-experts puissent comprendre et contester dans le cadre de
proc\'edures \'etablies~\cite{Daubert1993,Casey10}. Longtemps g\'erable, cette asym\'etrie s'est
creus\'ee avec les protocoles modernes au point de devenir une fracture structurelle~\cite{Citron14,Kerr04}.

Pour en d\'ecrire les contours, ce chapitre s'appuie sur trois syst\`emes
cryptographiques employ\'es comme construits illustratifs th\'eoriques, \`a savoir Helios pour le vote
\'electronique, les zk-SNARKs pour la conformit\'e financi\`ere et ECDSA pour
les signatures contractuelles. Leurs profils CLO ($\Conf$, $\Rel$, $\Opp$)
sont des estimations param\'etriques calcul\'ees analytiquement depuis leurs
sp\'ecifications publi\'ees, sous des hypoth\`eses institutionnelles d\'etaill\'ees
\`a la section~\ref{subsubsec:parametres_institutionnels}. Ils ne proviennent pas d'exp\'eriences
contr\^ol\'ees men\'ees avec des acteurs r\'eels. Ils servent \`a ancrer le paradoxe
dans des cas document\'es tout en motivant la formalisation des chapitres~\ref{chap:cadre_formel} \`a~\ref{chap:experimentation}.
\end{sloppypar}

% ----------------------------------------------------------------------------
\subsection{Cadre de recherche}
\label{subsec:cadre_recherche}
% ----------------------------------------------------------------------------
\begin{sloppypar}

\subsubsection{Le laboratoire LIMSI}
\label{subsubsec:limsi}

Les travaux pr\'esent\'es dans cette th\`ese ont \'et\'e conduits au
\textbf{Laboratoire d'Ing\'enierie Math\'ematique et Syst\`emes d'Information
(LIMSI)} de l'\'ecole Nationale Sup\'erieure Polytechnique de l'Universit\'e
de Yaound\'e~I, Cameroun. Le LIMSI accueille notamment des recherches situ\'ees \`a la
jonction de la cryptographie, des syst\`emes d'information, des m\'ethodes formelles et des math\'ematiques appliqu\'ees. Ce voisinage a permis de confronter les exigences de rigueur formelle aux difficult\'es que rencontre la pratique institutionnelle.

\subsubsection{L'exp\'erience professionnelle de l'auteur}
\label{subsubsec:experience_auteur}

Notre parcours professionnel associe plusieurs expertises compl\'ementaires, puisque nous exer\c{c}ons comme Auditeur des finances publiques au sein des \textbf{Services du Contr\^ole Sup\'erieur de l'\'Etat (CONSUPE)}, comme Expert judiciaire pr\`es les tribunaux du Cameroun ainsi que comme Expert certifi\'e en cybers\'ecurit\'e.

L'exercice de ces fonctions nous a souvent mis face \`a une situation d\'econcertante, dans laquelle des syst\`emes d'information robustes sur le plan cryptographique et s\^urs sur le plan math\'ematique ne parvenaient pas \`a convaincre les institutions qui devaient s'en servir~\cite{Bonneau16,Narayanan16}. En audit financier comme en expertise judiciaire, la validit\'e de la preuve technique suffisait rarement \`a fonder la confiance s\'emantique qu'exige une d\'ecision administrative ou juridictionnelle~\cite{Casey16,Grimmelmann15}.

\end{sloppypar}

% ----------------------------------------------------------------------------
\subsection{Gen\`ese de la question centrale de recherche}
\label{subsec:genese_question}
% ----------------------------------------------------------------------------
\begin{sloppypar}

La question centrale de cette th\`ese proc\`ede donc moins d'une r\'eflexion men\'ee \textit{in abstracto} que de l'observation r\'ep\'et\'ee d'un ph\'enom\`ene auquel nous avons donn\'e le nom de \textbf{paradoxe de l'authenticit\'e invisible}.

Le paradoxe tient \`a la conception m\^eme des protocoles cryptographiques actuels, qui ne prennent pas en compte la capacit\'e cognitive de leurs interpr\`etes~\cite{Simon82,Miller56,Herley09}, bien davantage qu'\`a un d\'efaut d'acculturation ou de formation des acteurs institutionnels.

Afin de d\'epasser le seul cadre national dans lequel ces constats avaient \'et\'e faits, nous avons retenu trois cas d'\'etude de port\'ee internationale (Helios pour le vote, les zk-SNARKs pour la conformit\'e, ECDSA pour la signature), dont les enseignements valent au-del\`a de leur contexte d'origine pour les questions de confiance num\'erique~\cite{DeFilippi16,Garay15}.

\end{sloppypar}

\begin{boiteArbitre}
L'arbitre a plusieurs dossiers sur sa table. Dans celui d'Helios figure une
\'election v\'erifiable de bout en bout que les garants \'electoraux n'ont pas su
d\'echiffrer sans assistance externe. Celui des zk-SNARKs contient une preuve de
quelques centaines d'octets qui atteste la conformit\'e de milliers de
transactions sans permettre au r\'egulateur de cibler la moindre op\'eration
suspecte. Le dossier d'ECDSA offre au contraire une transparence totale ainsi
qu'une tra\c{c}abilit\'e parfaite, obtenues au prix d'une confidentialit\'e nulle.
Chacun t\'emoigne d'une impuissance diff\'erente. Le pr\'esent chapitre les d\'ecrit
tandis que les chapitres suivants cherchent pourquoi elles convergent toutes
vers la m\^eme contrainte.
\end{boiteArbitre}

% ----------------------------------------------------------------------------
\subsection{Note de m\'ethodologie~: statut des exemples illustratifs}
\label{subsec:note_methodo}
% ----------------------------------------------------------------------------
\begin{sloppypar}
Cette th\`ese est une contribution th\'eorique. Les syst\`emes \'evoqu\'es au fil
du texte, Helios~\cite{Adida08,Adida09UCL}, ECDSA~\cite{Johnson01,NIST.FIPS.186-5},
zk-SNARKs~\cite{Groth16,Gennaro13,BCTV14} de conformit\'e, servent d'ancrages
illustratifs et de supports de validation th\'eorique, non de sujets d'\'etude
empirique. Les valeurs num\'eriques associ\'ees ($\Conf$, $\Rel$, $\Opp$) sont
des estimations param\'etriques d\'eriv\'ees analytiquement depuis les sp\'ecifications
publi\'ees de ces protocoles, sous des hypoth\`eses institutionnelles explicitement
formul\'ees (sous-section~\ref{subsubsec:hypotheses_capacite},
section~\ref{subsubsec:parametres_institutionnels}). Elles ne sont pas issues
d'exp\'eriences contr\^ol\'ees avec des acteurs institutionnels r\'eels. Ces
syst\`emes, dont la fiabilit\'e de v\'erification est proche de~$1$, ne sont pas
des instances du Trilemme CRO (Remarque~\ref{rem:portee_cas_etude}). Ils
servent \`a une \textbf{calibration illustrative} des ordres de grandeur des
param\`etres institutionnels, non \`a une validation, ni empirique au sens des
sciences exp\'erimentales~\cite{Tsybakov2009}, ni th\'eorique.
\end{sloppypar}

% ----------------------------------------------------------------------------
\subsection{Cas d'\'etude d\'eploy\'es}
\label{subsec:cas_etude}
% ----------------------------------------------------------------------------
\begin{sloppypar}
Les trois cas d'\'etude introduits \`a la section~\ref{subsec:genese_question} permettent de pr\'eciser le paradoxe. Leur choix tient \`a ce qu'ils
parcourent l'ensemble des compromis possibles entre confidentialit\'e et
opposabilit\'e, depuis les syst\`emes tr\`es confidentiels mais peu opposables
(Helios, zk-SNARKs) jusqu'\`a un syst\`eme tr\`es opposable dont la confidentialit\'e
est quasi nulle (ECDSA). Pris ensemble, ils font appara\^itre une r\'egularit\'e,
puisque aucun protocole d\'eploy\'e n'atteint un produit $\Conf \cdot \Rel \cdot \Opp$
sup\'erieur \`a $0{,}57$. Cette observation laisse supposer l'existence d'une
contrainte de fond.
\end{sloppypar}

% ----------------------------------------------------------------------------
\subsubsection{Vote \'electronique~: le cas Helios}
\label{subsubsec:cas_helios}
% ----------------------------------------------------------------------------
\begin{sloppypar}

\paragraph{Architecture cryptographique.}
\label{par:helios_architecture}

Le protocole Helios, con\c{c}u par Ben Adida en 2008~\cite{Adida08}, est
une r\'ef\'erence pour le vote \'electronique v\'erifiable de bout en bout~\cite{Benaloh94,Lindeman12}. Il
associe le chiffrement homomorphe ElGamal~\cite{ElGamal85}, gr\^ace auquel les
bulletins chiffr\'es peuvent \^etre agr\'eg\'es sans \^etre d\'echiffr\'es un \`a un, \`a des
preuves non-interactives zero-knowledge
(protocole de Chaum-Pedersen~\cite{ChaumPedersen92} rendu non interactif par Fiat-Shamir~\cite{FiatShamir87})
qui attestent la validit\'e de chaque bulletin sans en r\'ev\'eler le contenu.
Pour un bulletin binaire $m \in \{0,1\}$, le chiffr\'e est
$c = (g^r,\, h^r \cdot g^m)$ o\`u $h = g^x$ est la cl\'e publique
de l'\'election. Le d\'epouillement requiert la participation d'un seuil de garants ($t$
parmi $k$), la cl\'e de d\'echiffrement \'etant r\'epartie par le partage de secret de Shamir.

Les garanties qui en r\'esultent (confidentialit\'e individuelle, v\'erifiabilit\'e
universelle, r\'esistance partielle \`a la coercition~\cite{Benaloh94}) reposent
sur des fondements math\'ematiques solides~\cite{CramerDamgardSchoenmakers}. Rien
n'assure pour autant qu'elles soient accessibles aux institutions qui les
mettent en \oe uvre~\cite{Bonneau12,Whitten99}.

\paragraph{L'\'echec de Louvain (2009) et l'incident IACR (2025).}
\label{par:echec_louvain}

En mars 2009, l'Universit\'e Catholique de Louvain a d\'eploy\'e Helios pour
l'\'election de son Conseil \'etudiant~\cite{Adida09UCL}. Les garants
\'electoraux, des membres \'elus d\'epourvus de formation cryptographique, n'ont
pas pu conduire le d\'echiffrement distribu\'e sans assistance externe. La
v\'erifiabilit\'e universelle, pourtant garantie par construction, est ainsi
rest\'ee hors de port\'ee des observateurs auxquels elle \'etait destin\'ee.

En novembre 2025, l'IACR a d\^u annuler son \'election annuelle apr\`es qu'un
trustee eut perdu de fa\c{c}on irr\'em\'ediable sa cl\'e partielle~\cite{IACR2025incident}.
La configuration exigeait en effet les parts des trois garants (partage 3-sur-3).
Le scrutin a donc \'et\'e repris avec un seuil 2-sur-3. Les acteurs concern\'es
\textit{\'etaient} cette fois des experts en cryptographie, si bien que l'incident
met en \'evidence, ind\'ependamment de toute question de comp\'etence,
l'\'ecart entre la complexit\'e op\'erationnelle du protocole et les
contraintes concr\`etes d'un d\'eploiement institutionnel.

\paragraph{Profil CLO estim\'e.}
\label{par:helios_clo}

Sous les hypoth\`eses institutionnelles d\'etaill\'ees \`a la
section~\ref{sec:derivation_helios} et \`a la section~\ref{subsubsec:parametres_institutionnels},
les estimations param\'etriques analytiques retenues pour Helios sont
$\Conf \approx 0{,}92$ (confidentialit\'e op\'erationnelle du bulletin chiffr\'e ElGamal~\cite{TsiounisYung}, estim\'ee sur l'\'echelle CLO),
$\Relver \approx 0{,}95$ (compl\'etude et solidit\'e de la v\'erification des preuves de bien-formation),
$\Opp \approx 0{,}50$ (d\'ependance structurelle aux experts pour
l'interpr\'etation, avec capacit\'e d'extraction $\Delta \approx 56$ bits~\cite{MINKA2026CRO}),
ce qui donne le produit
\begin{equation}
\Conf \cdot \Rel \cdot \Opp \;\approx\; 0{,}437.
\label{eq:helios_cro}
\end{equation}
Cette valeur sert \`a la calibration illustrative des param\`etres
institutionnels (sous-section~\ref{subsec:analyse_comparative})~; le syst\`eme
n'est pas une instance du Trilemme CRO (Remarque~\ref{rem:portee_cas_etude}).

\end{sloppypar}

% ----------------------------------------------------------------------------
\subsubsection{Conformit\'e financi\`ere~: les zk-SNARKs}
\label{subsubsec:cas_zksnark}
% ----------------------------------------------------------------------------
\begin{sloppypar}

\paragraph{Succinctit\'e et opacit\'e.}
\label{par:zksnark_architecture}

Les arguments succincts non-interactifs zero-knowledge (zk-SNARKs)~\cite{Groth16,BCTV14,BenSasson2019}
ont nettement renforc\'e la confidentialit\'e des preuves. Une preuve $\pi$
de taille constante ($128$--$256$ octets) y certifie que
$\mathcal{C}(x, w) = 1$ pour un t\'emoin secret $w$, sa v\'erification ne
demandant que quelques millisecondes~\cite{Parno13,Wahby18}. Ces propri\'et\'es expliquent l'int\'er\^et
des zk-SNARKs pour la conformit\'e r\'eglementaire, puisqu'elles permettent de prouver
le respect des r\`egles AML~\cite{FATF12,AMLD4} ou de B\^ale~III~\cite{BaselIII}
sans divulguer les transactions sous-jacentes~\cite{Garratt19}.

Cette succinctit\'e a toutefois un co\^ut~\cite{Bitansky2017}, car une preuve qui ne r\'ev\`ele rien du t\'emoin
interdit aussi \`a un r\'egulateur d'auditer apr\`es coup une transaction suspecte.
Lors d'une crise syst\'emique comparable \`a celle des \textit{subprimes} de 2008~\cite{FCIC2011}, l'analyse de la contagion deviendrait impossible, l'opacit\'e
qui prot\`ege les clients faisant alors obstacle \`a l'investigation~\cite{Androulaki13,Meiklejohn13}. Le probl\`eme rel\`eve d'un vide de
contr\^ole~\cite{FinckMoscon19,DeFilippi16} bien plus que d'une d\'efaillance d'impl\'ementation.

\paragraph{Tensions avec le droit de la protection des donn\'ees.}
\label{par:zksnark_rgpd}

Le RGPD (R\`eglement G\'en\'eral pour la Protection des Donn\'ees \`a caract\`ere personnel)~\cite{GDPR2016} consacre un droit d'acc\`es (Article~15), un droit
\`a l'effacement (Article~17) ainsi qu'une obligation de rendre compte
(Article~5(2))~\cite{CCPA2018}. Chacune de ces exigences se heurte aux
propri\'et\'es zero-knowledge. Le droit d'acc\`es suppose en effet de pouvoir
d\'echiffrer, le droit \`a l'effacement s'oppose \`a l'immuabilit\'e de la
blockchain~\cite{Nakamoto08,Garay15} et l'obligation de justification pr\'esuppose
l'auditabilit\'e. La litt\'erature n'a pas encore r\'esolu ces
incompatibilit\'es~\cite{FinckMoscon19,Sergey18}.

\paragraph{Profil CLO estim\'e.}
\label{par:zksnark_clo}

Sous les m\^emes hypoth\`eses, les estimations param\'etriques analytiques donnent
$\Conf \approx 0{,}89$ (confidentialit\'e op\'erationnelle estim\'ee du dispositif de conformit\'e, au-del\`a de la seule preuve~\cite{Groth16}),
$\Relver \approx 0{,}96$ (compl\'etude et solidit\'e de la v\'erification),
$\Opp \approx 0{,}67$ (opposabilit\'e \'elev\'ee pour les institutions dot\'ees
d'une forte capacit\'e d'extraction $\Delta \approx 73$ bits~\cite{MINKA2026CRO}),
d'o\`u le produit
\begin{equation}
\Conf \cdot \Rel \cdot \Opp \;\approx\; 0{,}572.
\label{eq:zksnark_cro}
\end{equation}
Ce produit CRO relativement \'elev\'e s'explique par la forte capacit\'e
d'extraction des institutions financi\`eres sp\'ecialis\'ees ($\Delta \approx 73$ bits), qui disposent d'infrastructures
d'analyse, d'\'equipes d'experts et de cadres r\'eglementaires d\'etaill\'es~\cite{Ravi2020}.

\end{sloppypar}

% ----------------------------------------------------------------------------
\subsubsection{Signatures contractuelles~: ECDSA comme r\'ef\'erence d'opposabilit\'e}
\label{subsubsec:cas_ecdsa}
% ----------------------------------------------------------------------------
\begin{sloppypar}

\paragraph{Architecture et acceptation universelle.}
\label{par:ecdsa_architecture}

L'\emph{Elliptic Curve Digital Signature Algorithm} (ECDSA)~\cite{Johnson01}
se situe \`a l'autre extr\'emit\'e du spectre. Normalis\'e par le NIST dans FIPS~186-2 (2000) puis FIPS~186-5 (2023)~\cite{NIST.FIPS.186-5},
il est d\'eploy\'e dans SSL/TLS, dans Bitcoin~\cite{Nakamoto08,Narayanan16} ainsi que dans
les syst\`emes bancaires. Sa s\'ecurit\'e repose sur la difficult\'e du probl\`eme du logarithme discret sur
courbes elliptiques (ECDLP)~\cite{RFC6090,RFC6979}, la cl\'e publique $Q = d \cdot G$
\'etant obtenue \`a partir de la cl\'e secr\`ete $d \in [1, n-1]$. La signature d'un message $m$ produit
$(r, s)$ v\'erifiable par quiconque conna\^it $Q$ et $m$.

Cette transparence totale fonde l'opposabilit\'e d'ECDSA~\cite{ESIGN00,eIDAS14}, puisque la
v\'erification ne d\'epend d'aucun expert tiers, peut \^etre rejou\'ee \`a tout moment et
aboutit \`a une attribution univoque assortie d'une tra\c{c}abilit\'e compl\`ete~\cite{Sasson14}.
Elle a pour contrepartie directe l'absence de toute confidentialit\'e, le message
circulant en clair tandis que l'identit\'e est publique et la transaction trac\'ee~\cite{Androulaki13,ECDSA-Fails}.

\paragraph{Profil CLO estim\'e.}
\label{par:ecdsa_clo}

Les estimations param\'etriques analytiques sont
$\Conf \approx 0{,}31$ (confidentialit\'e op\'erationnelle estim\'ee du dispositif contractuel~; au sens statistique de la D\'efinition~\ref{def:confidentialite_clo}, un message publi\'e en clair a une confidentialit\'e nulle),
$\Relver \approx 0{,}98$ (infalsifiabilit\'e existentielle sous ECDLP~\cite{GoldwasserMicali84,GMR88}),
$\Opp \approx 0{,}90$ (v\'erification rejouable, non-r\'epudiation forte,
capacit\'e d'extraction $\Delta \approx 35$ bits suffisante~\cite{MINKA2026CRO}),
soit le produit
\begin{equation}
\Conf \cdot \Rel \cdot \Opp \;\approx\; 0{,}273.
\label{eq:ecdsa_cro}
\end{equation}
La faiblesse de ce produit ne signale aucune d\'efaillance de conception.
Elle d\'ecoule d'un choix architectural d\'elib\'er\'e, ECDSA privil\'egiant
l'opposabilit\'e au d\'etriment de la confidentialit\'e~\cite{BonehFranklin}.

\end{sloppypar}

% ----------------------------------------------------------------------------
\subsubsection{Analyse comparative et lacune th\'eorique}
\label{subsubsec:analyse_comparative}
% ----------------------------------------------------------------------------
\begin{sloppypar}
Le Tableau~\ref{tab:clo_resume_ch1} r\'esume les trois profils pr\'esent\'es plus haut
et appelle deux observations. En premier lieu, confidentialit\'e et opposabilit\'e
varient en sens inverse, les protocoles qui maximisent l'une minimisant
l'autre~\cite{ClarksonSchneider10}. Aucun des trois n'approche la r\'egion $(\Conf, \Rel, \Opp)
\approx (0{,}8, 0{,}8, 0{,}8)$, qui resterait pourtant en dessous de
$1$. Ensuite, les produits $\Conf \cdot \Rel \cdot \Opp$ varient de
$0{,}273$ \`a $0{,}572$, ce qui sugg\`ere l'existence d'une contrainte
qui d\'epend des param\`etres
institutionnels ($\Delta$, $\alpha$, $\beta$)~\cite{Simon82,Cowan2001}.

La cryptographie classique ne dispose d'aucun outil pour formaliser cette
contrainte~\cite{AbadiRogaway00}, car elle mesure l'avantage de l'adversaire~\cite{Goldreich01} sans rien dire de l'utilit\'e institutionnelle. La litt\'erature sur l'\emph{usable
security}~\cite{Bonneau12,Cranor05,Herley09,Dhamija06,Schechter07} \'etudie l'utilisabilit\'e, mais
qualitativement, sans m\'etriques quantifiables. L'informatique
juridique~\cite{Kerr04,Casey11,Mason07} analyse l'admissibilit\'e des preuves,
mais via la jurisprudence, sans formalisme math\'ematique.

Aucune de ces trois disciplines, qui s'ignorent largement, ne propose donc de
m\'etrique de l'opposabilit\'e~\cite{Smith09,Alvim12}. La section~\ref{subsec:trois_dimensions} s'attache \`a combler ce manque.
\end{sloppypar}

\begin{table}[htbp]
\centering
\caption{Profils CLO estimés~: trois systèmes représentatifs.
Toutes les valeurs sont des estimations paramétriques analytiques publiées
dans les travaux CLO~\cite{MINKA2026CLO} et CRO~\cite{MINKA2026CRO}, sur l'\'echelle
op\'erationnelle CLO (fiabilit\'e de v\'erification $\Relver$, Remarque~\ref{rem:rel_ver}).}
\label{tab:clo_resume_ch1}
\begin{tabular}{lcccc}
\toprule
\textbf{Système} & $\Conf$ & $\Relver$ & $\Opp$ & $\Conf \cdot \Relver \cdot \Opp$ \\
\midrule
Helios (vote)        & $0{,}92$ & $0{,}95$ & $0{,}50$ & $0{,}437$ \\
zk-SNARKs (finance) & $0{,}89$ & $0{,}96$ & $0{,}67$ & $0{,}572$ \\
ECDSA (contrats)    & $0{,}31$ & $0{,}98$ & $0{,}90$ & $0{,}273$ \\
\bottomrule
\end{tabular}
\end{table}

% ----------------------------------------------------------------------------
\subsection{Extension quantique~: les standards NIST PQC}
\label{subsec:extension_quantique_contexte}
% ----------------------------------------------------------------------------
\begin{sloppypar}

\subsubsection{La transition quantique}
\label{subsubsec:transition_quantique}

En ao\^ut 2024, le National Institute of Standards and Technology (NIST) a publi\'e
les trois premi\`eres normes de cryptographie post-quantique~\cite{NISTPQC,Alagic2020,NISTPostQuantum2022}, qui ouvrent une nouvelle
p\'eriode pour la discipline~\cite{Chen2022}. Il s'agit des normes suivantes.

\begin{itemize}
\item \textbf{FIPS~203 (ML-KEM)}, m\'ecanisme d'encapsulation de cl\'es fond\'e sur
les r\'eseaux euclidiens et d\'eriv\'e de Kyber~\cite{Kyber2022}, sert \`a l'\'etablissement de cl\'es
ainsi qu'au chiffrement hybride.

\item \textbf{FIPS~204 (ML-DSA)}, algorithme de signature num\'erique fond\'e sur
les r\'eseaux euclidiens et d\'eriv\'e de Dilithium~\cite{Dilithium2022}, est appel\'e \`a remplacer ECDSA et RSA
pour les signatures.

\item \textbf{FIPS~205 (SLH-DSA)}, algorithme de signature fond\'e sur les fonctions
de hachage et d\'eriv\'e de SPHINCS+~\cite{SPHINCS2022}, constitue une alternative sans trusted setup qui repose sur
des hypoth\`eses cryptographiques conservatrices.
\end{itemize}

Le standard FIPS~206 (FN-DSA), d\'eriv\'e de Falcon~\cite{Falcon2022}, en cours de finalisation, compl\'etera cet ensemble.

\subsubsection{Enjeux institutionnels de la migration PQC}
\label{subsubsec:enjeux_pqc}

La migration vers la cryptographie post-quantique fait na\^itre de nouvelles
tensions avec les exigences institutionnelles~\cite{Langlois2015,Boneh2011}, qui tiennent \`a la complexit\'e des preuves, \`a leur taille et \`a la nouveaut\'e des hypoth\`eses de s\'ecurit\'e.

\paragraph{Complexit\'e accrue des preuves.}
Les signatures Dilithium reposent sur des preuves de connaissance dans les r\'eseaux
euclidiens, dont la v\'erification fait intervenir des op\'erations matricielles et des
calculs de normes~\cite{Dilithium2022}. Un arbitre habitu\'e \`a ECDSA devra donc se familiariser avec les
modules sur $\mathbb{Z}_q$ ainsi qu'avec les probl\`emes LWE (Learning With Errors)~\cite{Regev2009}.

\paragraph{Taille des preuves.}
Une signature ML-DSA-65 (Dilithium-3) occupe 3\,309 octets et sa cl\'e publique 1\,952 octets, contre 64 octets pour une signature ECDSA-256.
De m\^eme, une cl\'e publique Kyber-1024 atteint 1\,568 octets, l\`a o\`u ECDSA se contente de 64 octets.
Cet accroissement complique l'examen manuel des preuves comme leur archivage~\cite{Ravi2020}.

\paragraph{Hypoth\`eses de s\'ecurit\'e non standard.}
LWE et les probl\`emes de r\'eseaux euclidiens constituent des hypoth\`eses r\'ecentes
(ann\'ees 2000), qui ne b\'en\'eficient pas de la confiance accumul\'ee au fil du temps par RSA (1977) ou
ECDLP (1985)~\cite{Kiltz2017}. Des institutions prudentes pourraient donc h\'esiter \`a reconna\^itre la m\^eme
force probante \`a des preuves fond\'ees sur ces hypoth\`eses nouvelles~\cite{Unruh2016}.

\subsubsection{Implications pour le Trilemme CRO}
\label{subsubsec:implications_cro}

La transition post-quantique conduit \`a se demander si les bornes du Trilemme CRO
\'etablies dans le contexte classique restent valides dans les trois situations
suivantes~\cite{NielsenChuang2010,Wilde2017}.

\begin{enumerate}
\item L'adversaire peut effectuer des requ\^etes quantiques en superposition~\cite{Watrous2009,Boneh2011}.
\item Les canaux de communication introduisent des d\'egradations quantiques~\cite{Holevo1973,Schumacher1997}.
\item Les param\`etres institutionnels ($\Delta$, $\alpha$, $\beta$) doivent
\^etre adapt\'es au contexte quantique~\cite{Mayers1997,Lo1997}.
\end{enumerate}

C'est pour y r\'epondre que la section~\ref{subsec:modele_formel_pq} d\'eveloppe
l'extension quantique du Trilemme CRO, en introduisant les param\`etres
$\Delta_q$ (capacit\'e d'extraction quantique) et $\eta_q$ (perte d'interpr\'etabilit\'e
quantique).

\end{sloppypar}

% ----------------------------------------------------------------------------
%\subsection{Synth\`ese}
%\label{sec:synthese_chap1}
% ----------------------------------------------------------------------------

\begin{boiteArbitre}
L'arbitre conna\^it ces dossiers pour les avoir tenus entre ses mains.
Helios, les zk-SNARKs et ECDSA reposent sur des architectures distinctes,
mais chacun a \'echou\'e \`a sa mani\`ere \`a le convaincre.

Ce qui l'a longtemps troubl\'e tient \`a ce que la d\'efaillance n'\'etait jamais
d'ordre technique. Les preuves \'etaient correctes et les algorithmes
irr\'eprochables, sans que cela lui permette pour autant de d\'ecider. Il lui
manquait moins la comp\'etence qu'un langage commun entre la rigueur
math\'ematique du protocole et les exigences proc\'edurales de sa fonction. Ce
chapitre a donn\'e \`a cette situation le nom de paradoxe de l'authenticit\'e
invisible.

Il en retient quatre enseignements. Une confidentialit\'e tr\`es \'elev\'ee
(sup\'erieure \`a $0{,}90$) produit une opacit\'e qui paralyse l'interpr\`ete
institutionnel, comme le montrent Helios et les zk-SNARKs. Les
produits $\Conf \cdot \Rel \cdot \Opp$ s'\'etagent de $0{,}273$ \`a $0{,}572$
selon les param\`etres institutionnels ($\Delta$, $\alpha$, $\beta$), ce qui
d\'esigne une contrainte qui d\'epend du contexte.
Les protocoles actuels ont aussi \'et\'e con\c{c}us sans tenir compte de
lui, puisque sa capacit\'e cognitive born\'ee ($\Delta$) n'appara\^it dans aucune
sp\'ecification. Enfin, la transition post-quantique risque d'aggraver cette
situation en y ajoutant une complexit\'e que les experts eux-m\^emes peinent \`a
ma\^itriser.

Il reste \`a passer du constat au th\'eor\`eme, ce que le chapitre suivant
pr\'epare en posant l'infrastructure th\'eorique n\'ecessaire.
\end{boiteArbitre}